%#! latex main.tex
\appendix
\addtocontents{toc}{\protect\par\vskip10mm}

\chapter{Rapid Course in $p$-adic Analysis.}\label{app:rapid}
\renewcommand{\thechapter}{\Alph{chapter}}
\addtocontents{toc}{\protect\par\vskip2mm\hskip80mm by F. Kato \par\vskip5mm}

\vspace{1ex}
\begin{center}
By {\sc F.\ Kato}
\end{center}

%
\section{Introduction.}
%
%Let us first set up the notation; in order to do it, and for the
%reader's convenience, we recall some terminologies:
%\begin{ter}
%A {\it semi-norm} on a commutative and unitary ring $A$ is a 
%function $\|\cdot\|\colon A\rightarrow\R_{\geq 0}$ which satisfies
%$\|0\|=0$, $\|1\|=1$, $\|fg\|\leq\|f\|\|g\|$, and $\|f-g\|\leq\|f\|+
%\|g\|$ for $f,g\in A$.
%We often write $\Ker\|\cdot\|\colon=\{f\in A\ |\ \|f\|=0\}$.
%In case $\Ker\|\cdot\|=\{0\}$ we say $\|\cdot\|$ is a {\it norm}.
%A semi-norm $\|\cdot\|$ on a ring $A$ is said to be {\it non-archimedan}
%if for $f,g\in A$ $\|f-g\|\leq\mathrm{max}\{\|f\|,\|g\|\}$.
%A norm $\|\cdot\|$ is said to be a {\it valuation} if for $f,g\in A$
%$\|fg\|=\|f\|\|g\|$.
%For example, the function $\|\cdot\|_0$ which sends every non-zero element
%in $A$ to $1$ is a norm; this norm is called the {\it trivial norm}.
%\end{ter}

In this appendix, $K$ always denotes a complete field 
with respect to a non-archimedean valuation $|\cdot|\colon K\rightarrow
\R_{\geq 0}$.
The norm $|\cdot|$ is almost always assumed to be non-trivial, unless
otherwise stated.

Let us assume for a while that $K$ is algebraically closed.
We view the field $K$ as an affine space, as we do in complex analysis; 
$K$ is a metrized space with the metric induced from the
valuation $|\cdot|$.
We can then follow the lines of classical complex analysis and define
convergence of power series, Taylor and Laurent
expansions, etc...
%to make sense in this situation.
%Then we easily find that it is quite meaningful for us to {\it ask} 
%if a given $K$-valued
%function defined on an open set in $K$ is ``analytic'', i.e.\
%locally represented by a convergent power series or not.
%So we are tempted to {\it characterize} those analytic or holomorphic 
%functions by the last condition, also as we did in complex analysis; but 
%this is in fact very problematic.
%
%To see what happens more clearly, let us test this na\"{\i}ve approach; 

For instance let us define ``holomorphic'' functions to be $K$-valued
functions which are locally expressed by convergent power
series\footnote{The reader may wonder whether there is another approach
by means of Cauchy-Riemann type differential equations.  This is unlikely
because the differential calculus does not go well, see
\cite[Chap.\ 4]{pn}.}  (cf.\ \cite[Chap.\ 4]{pn}). 
This approach drives us however to several
problems, which come mainly from the fact that $K$ with the
metric topology is totally disconnected; for example:
\begin{enumerate}
 \item for an open set $U\subseteq K$ the ring of all such holomorphic
       functions on $U$ is huge. So is already the subring of all locally 
       constant functions.
 \item It can be shown that the sheaf of germs of such functions satisfies
       the principle of unique continuation; but not in a satisfactory way.
       In particular, \lemref{lem-unique-cont} of \chapref{chap-ana-p-period}
       holds but is of no use,
       because there is no non-empty connected open subset.
%\item There is no reasonable path; any continuous map $\gamma\colon
%[0,1]\rightarrow K$ is constant.
\end{enumerate}

The trouble becomes more apparent when we think the \emph{local}
representability by convergent power series in terms of
\emph{coverings}.  In the complex analytic situation, when we speak of
the holomorphy of a function $f$ defined over a connected open set $U$,
we tacitly take a open covering $\{U_i\}_{i\in I}$ of $U$ consisting of
sufficiently small open neighborhoods such that $f$ restricted on each
member $U_i$ can be seen as evaluation of a power series convergent in
$U_i$.  The point is that, since $U$ is connected, members of $\{U_i\}$
must have overlaps so that the local properties can be transmitted to
whole $U$.  But, in the non-archimedean situation, such coverings may be
refined in such a way that there are no overlap.  This is why
the analytic continuation does not work well.  So, for the sake of a
reasonable analytic theory, we have to ``limit'' coverings so that we
cannot take arbitrarily fine refinements.

A first method is due to Krasner.
He introduced the so-called {\it quasi-connected sets} and, on such
subsets, defined analytic functions as uniform convergent limits of
rational functions.

A more modern and systematic treatment was introduced by Tate, 
with the so-called {\it Rigid Analysis}.
Let us briefly view the main idea:
we first set
%$$
%K\{t_1,\ldots,t_n\}=\left\{
%\begin{array}{r|l}
%\sum_{\nu_1,\ldots,\nu_n\geq 0}
%a_{\nu_1,\ldots,\nu_n}t_1^{\nu_1}\cdots t_n^{\nu_n}&
%|a_{\nu_1,\ldots,\nu_n}|\rightarrow 0\ \mathrm{for}\\
%\in K[[t_1,\ldots,t_n]]
%&\nu_1+\cdots+\nu_n\rightarrow\infty
%\end{array}
%\right\},
%$$
\begin{align*}
   & K\{t_1,\ldots,t_n\}\\
 = & \Biggl\{
 \sum_{\nu_1,\ldots,\nu_n\geq 0}
 a_{\nu_1,\ldots,\nu_n}t_1^{\nu_1}\cdots t_n^{\nu_n}
  \in K[[t_1,\ldots,t_n]]\,
 \Bigg|
 \begin{split}
  & a_{\nu_1,\ldots,\nu_n}\rightarrow 0\ \mathrm{for}\\
  & \nu_1+\cdots+\nu_n\rightarrow\infty
 \end{split} 
\Biggr\},
\end{align*}
which is now called the {\it Tate algebra} over $K$.
This is the ring of all functions expressed by power series convergent
on the closed disk $\matheur{D}(0,1^{+})=\{z\in K^n\ |\ |z_i|\leq 1\}$.
The algebra $K\{t_1,\ldots,t_n\}$ is endowed with the sup-norm 
$\|\cdot\|$, also called the {\it Gauss norm}.
We list up some known properties (cf.\ \cite[5.2]{na},
\cite[II.3]{garea}):

\begin{enumerate}
\item The ring $K\{t_1,\ldots,t_n\}$ satisfies the Weierstrass
Preparation Theorem.
\item The ring $K\{t_1,\ldots,t_n\}$ is Noetherian and factorial.
\item Every ideal of $K\{t_1,\ldots,t_n\}$ is closed.
\item For any maximal ideal $\mathfrak{m}$ of $K\{t_1,\ldots,t_n\}$, the
field $K\{t_1,\ldots,t_n\}/\mathfrak{m}$ is a finite extension of $K$.
\end{enumerate}

If, moreover, $K$ is algebraically closed, the set of all maximal
ideals, endowed with the natural topology, coincides with 
$\matheur{D}(0,1^{+})$.
So the situation is analogous to classical algebraic geometry; 
$\matheur{D}(0,1^{+})$ to the affine $n$-space, and $K\{t_1,\ldots,t_n\}$
to the coordinate ring.
In fact, quotients $\matheur{A}=K\{t_1,\ldots,t_n\}/I$ by ideals $I$ 
(called {\it affinoid algebras}), and their associated maximal spectra,
$\mathrm{Spm}(\matheur{A}$), called {\it affinoids}, forms the 
fundamental patches from which the non-archimedean function theory will be
developed.
Although the natural metric topology on $\mathrm{Spm}(\matheur{A})$ is 
terrible for the same reason as above, the situation becomes much
better as far as we deal only with those functions coming from 
$\matheur{A}$.
The trick is that, if the algebra $\matheur{A}$ is integral, then we can
pretend that the space $\mathrm{Spm}(\matheur{A})$ is ``connected''; for
example, the closed disk $\matheur{D}(0,1^{+})$ is ``connected'' in this
sense.
So, basically according to this point of view, we can globalize 
the situation by gluing affinoid patches to obtain a reasonable
theory of analytic functions.
The actual recipe to do it is furnished by the notion of Grothendieck
topology, which censors the plethora of coverings.
We will see this a little more precisely in what follows.

\begin{apprem}
Although the rigid analysis seems to provide the reasonable topological
and analytical framework, one must not expect it to be as nice as the
complex case.  Firstly, in classical complex analysis, when we expand
a holomorphic function centered at a point inside its region of
convergence, the resulting Taylor expansion may have a different region
of convergence than the original one. This is a important method of
analytic continuation.  But in the non-archimedean case,
this never happens (cf.\ \cite[4.4]{pn}).

Secondly, the topological framework
given as above is still not powerful enough to handle paths.
So it is hopeless to treat, for example, fundamental groups and 
monodromy in an intuitive way as in the classical complex case.
This difficulty is, in fact, remedied by Berkovich' theory 
which we shall see later on.
\end{apprem}

%
\section{Rigid analytic spaces.}
%

Let $K$ be as in the previous section; we do not assume in general 
that $K$ is algebraically closed.

As mentioned above, an {\it affinoid algebra}
$\matheur{A}$ over $K$ is a quotient $K\{t_1,\ldots,t_n\}/I$ for some
$n$ by an ideal $I$.  Since $I$ is closed, $\matheur{A}$ has a norm
$|\cdot |$ induced from the sup-norm $\|\cdot \|$ of
$K\{t_1,\ldots,t_n\}$ ({\it ``residue norm''}).  With the residue norm,
$\matheur{A}$ is a Banach $K$-algebra.  The norm $|\cdot |$ itself
depends on the presentation $\alpha\colon
K\{t_1,\ldots,t_n\}/I\stackrel{\sim}{\rightarrow}\matheur{A}$, while the
induced topology does not.  So, strictly speaking, we should write it
like $|\cdot |_{\alpha}$.  Clearly, $\matheur{A}$ is Noetherian and, for
any maximal ideal $\mathfrak{m}$ of $\matheur{A}$, the residue field
$\matheur{A}/\mathfrak{m}$ is a finite extension of $K$, and hence, the
valuation $|\cdot |$ of $K$ naturally extends to that of it, denoted
again by $|\cdot |$.

The associated {\it affinoid} $\mathrm{Spm}({\matheur{A}})$ is the set of all
maximal ideals of $\matheur{A}$.
For $x\in\mathrm{Spm}(\matheur{A})$ and $f\in\matheur{A}$ the value of $f$ at
$x$, denoted by $f(x)$, is the class of $f$ in $\matheur{A}/x$.
The set $\mathrm{Spm}({\matheur{A}})$ has the topology generated
by the subsets of form $\{x\in\mathrm{Spm}({\matheur{A}})\ |\ |f(x)|
\leq 1\}$ for $f\in\matheur{A}$.
But, as we pointed out above, this topology is not very interesting since 
it makes $\mathrm{Spm}({\matheur{A}})$ totally disconnected.
So we should specify the reasonable family of ``admissible'' open
sets (and coverings), on which the function theory will be built.

Let $\matheur{A}$ be an affinoid algebra over $K$, and $f_i$ 
($i=0,\ldots,n$) a collection of elements in $\matheur{A}$ which have no
common zeros on $\mathrm{Spm}(\matheur{A})$. 
The subspace
\begin{equation*}
 R=\{x\in\mathrm{Spm}(\matheur{A})\ |\ |f_i(x)|\leq |f_0(x)|,\ i=1,\ldots,n\}
\end{equation*}
can be identified with the affinoid $\mathrm{Spm}(\matheur{A}_R)$, where
\begin{equation*}
 \matheur{A}_R=\matheur{A}\widehat{\otimes}_KK\{t_1,\ldots,t_n\}/
 (f_1-t_1f_0,\ldots,f_n-t_nf_0).
\end{equation*}
A subset of this form is called a {\it rational subdomain}.
The identification comes as follows:
We first note that the morphism of affinoids $\mathrm{Spm}(\matheur{A}_R)
\rightarrow\mathrm{Spm}(\matheur{A})$ induced by
$\matheur{A}\rightarrow\matheur{A}_R$ maps $\mathrm{Spm}(\matheur{A}_R)$ to 
$R$.
Then $\matheur{A}_R$ is the unique solution of the following 
universal property (hence, it can be determined up to canonical isomorphism):
for any morphism of affinoids
$\phi\colon\mathrm{Spm}(\matheur{B})\rightarrow\mathrm{Spm}(\matheur{A})$ such
that $\phi(\mathrm{Spm}(\matheur{B}))\subset R$, there exists a unique 
$K$-homomorphism $\matheur{A}_R\rightarrow\matheur{B}$ such that
the resulting diagram
\vspace{1.5ex}
$$
\setlength{\unitlength}{1pt}
\begin{picture}(100,60)(10,0)
\put(0,45){$\mathrm{Spm}(\matheur{B})$}
\put(40,48){\vector(1,0){20}}
\put(46,53){$\scriptstyle{\phi}$}
\put(66,45){$\mathrm{Spm}(\matheur{A})$}
\put(29,6){$\mathrm{Spm}(\matheur{A}_R)$}
%\put(33,38){\vector(1,-2){10}}
\multiput(33,38)(1.5,-3){6}{\circle*{.3}}\put(42,20){\vector(1,-2){2}}
\put(57,18){\vector(1,2){10}}
\end{picture}\vspace{-1ex}
$$
is commutative.

\begin{applem}\label{lem-rational-intersection}

$\mathrm{(1)}$
If $R$ and $S$ are rational subdomains in $\mathrm{Spm}(\matheur{A})$, then 
so is $R\cap S$.

$\mathrm{(2)}$
If $R_1$ is a rational subdomain in $\mathrm{Spm}(\matheur{A})$, and $R_2$ 
is a rational subdomain in $R_1$, then $R_2$ is a rational subdomain 
in $\mathrm{Spm}(\matheur{A})$.
\end{applem}

\begin{proof}
(1)
If $R=\{|f_i(x)|\leq |f_0(x)|,\ i=1,\ldots,n\}$ and
$S=\{|g_j(x)|\leq |g_0(x)|,\ j=1,\ldots,m\}$, then
one sees easily that 
$R\cap S=\{|f_i(x)g_j(x)|\leq |f_0(x)g_0(x)|,\ i=1,\ldots,n,
\ j=1,\ldots,m\}$.

(2) Exercise.
\end{proof}

\begin{appcor}\label{cor-rational-intersection}
Any subspace of the form
\begin{equation*}
 \{|f_i(x)|\leq 1,|g_j(x)|\geq 1,\ i=1,\ldots,n,\ j=1,\ldots,m\}
\end{equation*}
is a rational subdomain.
\end{appcor}

\begin{proof}
It is the intersection of rational subdomains
$\{|f_i(x)|\leq 1\}$ for $i=1,\ldots,n$ and
$\{1\leq |g_j(x)|\}$ for $j=1,\ldots,m$.
\end{proof}

\begin{appexa}[Rational subdomains of the unit polydisk]
Let us assume that $K$ is algebraically closed.
We consider rational subdomains of the unit polydisk
$\matheur{D}(0,1^{+})=\mathrm{Spm}(K\{t_1,\ldots,t_n\})$.

(1) {\it Closed polydisk}:\  
$\matheur{D}(0,|\pi|^{+})=\{x\in K^n\ |\ |x_i|\leq |\pi_i|,\ 
i=1,\ldots,n\}$,
for $0\neq\pi_i\in K$, $|\pi_i|\leq 1$, is a rational subdomain. 
The corresponding affinoid algebra is 
$$
K\left\{\frac{t_1}{\pi_1},\ldots,\frac{t_n}{\pi_n}\right\}.
$$

(2) {\it Annulus}:\ 
$\matheur{C}(0,|\pi^{(1)}|^{+},|\pi^{(2)}|^{+})=
\{x\in K^n\ |\ |\pi^{(1)}_i|\leq |x_i|\leq
|\pi^{(2)}_i|,\ i=1,\ldots,n\}$, 
for $0\neq\pi^{(j)}_i\in K$, $|\pi^{(j)}_i|\leq 1$, is a rational 
subdomain. The corresponding affinoid algebra is 
$$
K\left\{\frac{\pi^{(1)}_1}{t_1},\ldots,\frac{\pi^{(1)}_n}{t_n},
\frac{t_1}{\pi^{(2)}_1},\ldots,\frac{t_n}{\pi^{(2)}_n}\right\}.
$$
\end{appexa}

Let us next review the definition of Grothendieck topology
(in a narrow sense). Let $X$ be a topological space.
A {\it Grothendieck topology} (G-topology, in short) on $X$ is a pair 
$(\mathscr{T},\mathscr{C}\matheur{ov})$ consisting of
\begin{itemize}
\item a collection $\mathscr{T}$ of open subsets in $X$,
\item an assignment $U\mapsto\mathscr{C}\matheur{ov}(U)$ for any $U\in\mathscr{T}$,
where $\mathscr{C}\matheur{ov}(U)$ is a collection of coverings by elements in
$\mathscr{T}$,
\end{itemize}

\noindent
such that the following conditions are satisfied:

\vspace{1ex}
(1) $\emptyset\in\mathscr{T}$; $U,V\in\mathscr{T}\ \Rightarrow\ 
U\cap V\in\mathscr{T}$.

(2) $U\in\mathscr{T}\ \Rightarrow\ \{U\}\in\mathscr{C}\matheur{ov}(U)$.

(3) $\{U_i\}_{i\in I}\in\mathscr{C}\matheur{ov}(U),\ V\subseteq U,\ V\in
\mathscr{T}\ \Rightarrow\ \{U_i\cap V\}_{i\in I}\in
\mathscr{C}\matheur{ov}(V)$.

(4) $\{U_i\}_{i\in I}\in\mathscr{C}\matheur{ov}(U),\ \{U_{i,j}\}_{j\in J_i}
\in\mathscr{C}\matheur{ov}(U_i)\ \Rightarrow\ \{U_{i,j}\}_{i\in I,j\in J}
\in\mathscr{C}\matheur{ov}(U)$.

\vspace{1ex}
Elements in $\mathscr{T}$ are called {\it admissible open sets}, and elements
in $\mathscr{C}\matheur{ov}(U)$ are called {\it admissible coverings} of $U$.

Let $X=\mathrm{Spm}(\matheur{A})$ be an affinoid over $K$.
We introduce a Grothendieck topology to $X$
by the following recipe:

\begin{itemize}
\item An admissible open set is a rational subdomain.
\item For a rational subdomain, an admissible covering is a finite 
covering consisting of rational subdomains.
\end{itemize}
This is valid, due to Lemma \ref{lem-rational-intersection}.

We define the structure presheaf $\OO_X$, with respect to the Grothendieck
topology, by assigning for each rational subdomain $R$ the corresponding
affinoid algebra $\matheur{A}_R$.
Due to the following weak form of {\it Tate's acyclicity theorem}, the
presheaf $\OO_X$ is, in fact, a sheaf (a sheaf of local rings).

\begin{appthm}
Let $U_1,\ldots,U_m$ be rational subdomains of
$X=\mathrm{Spm}(\matheur{A})$. Set $U=U_1\cup\cdots\cup U_m$.
Then the sequence
$$
0\longrightarrow\OO_X(U)\longrightarrow\prod^m_{i=1}\OO_X(U_i)
\longrightarrow\prod^m_{i,j=1}\OO_X(U_i\cap U_j)
$$
is exact, where the last arrow is the difference of the two possible
restriction morphisms.
\end{appthm}

For the proof, see, for example, \cite[8.2]{na} or \cite[III.2.2]{garea}.

\begin{appdfn}[Rigid Analytic Space]
A {\it rigid analytic space} is a locally ringed space $(X,\OO_X)$, with 
a Grothendieck topology in the above sense, locally isomorphic to an
affinoid;
more precisely, there exists a covering $\{X_i\}_{i\in I}$ 
(possibly infinite) of $X$ by 
admissible open sets such that $(X_i,\OO_X|_{X_i})$ is isomorphic to 
a certain affinoid for each $i$.
\end{appdfn}

Let us see some examples of rigid analytic spaces:

\begin{appexa}[Projective Space]
The projective space $\mathbb{P}^n(K)$ has the natural rigid 
analytic structure.
Here we limit ourselves to demonstrate it only in the case $n=1$, and 
leave the general case to the reader.
Let $(X\colon Y)$ be the homogeneous coordinate in $\mathbb{P}^1(K)$
and set $z=X/Y$.
Set 
$$
\matheur{U}^{+}=\left\{(X\colon Y)\ |\ |X|\leq |Y|\right\}
\ \mathrm{and}\ 
\matheur{U}^{-}=\left\{(X\colon Y)\ |\ |X|\geq |Y|\right\}.
$$
Then $\matheur{U}^{+}$ is isomorphic to the closed disk
$\{z\in K\ |\ |z|\leq 1\}$ (the Southern Hemisphere with the equator)
, and $\matheur{U}^{-}$ to
$\{z\in K\cup\{\infty\}\ |\ |z|\geq 1\}$
(the Northern Hemisphere with the equator).
The intersection $\matheur{U}^{+}\cap\matheur{U}^{-}$ is thus isomorphic to
the circle $\{z\in K\ |\ |z|=1\}$.
Hence $\{\matheur{U}^{+},\matheur{U}^{-}\}$ gives the admissible covering of
$\mathbb{P}^1(K)$ which induces, in the obvious way, the structure
of a rigid analytic space.
\end{appexa}

More generally, an algebraic variety, separated and of finite type
over $K$, carries the canonical structure of a rigid analytic space.
%(cf.\ \cite[III.4.6]{garea}).

Any open set (not necessarily admissible) $U$ in a rigid analytic 
space $X$ will be again a rigid analytic space ({\it ``open subanalytic
space''}); more precisely, $U$ has the induced analytic structure from 
that of $X$. The admissible open sets of $U$, for example, are those
of $X$ contained in $U$.

%
\section{Relation with Formal Geometry.}
%

In this section we need to consider the valuation ring 
$R$ of $K$. Let us denote by $\mathfrak{m}$ the maximal ideal of $R$, and
fix a topological generator $0\neq\pi\in\mathfrak{m}$.
The residue field $R/\mathfrak{m}$ is denoted by $k$.

Let $\matheur{A}$ be an affinoid algebra over $K$.
The {\it spectral semi-norm} on $\matheur{A}$ is the function
$|\cdot |_{\mathrm{sup}}\colon\matheur{A}\rightarrow\R_{\geq 0}$
defined, for any $f\in\matheur{A}$, by
\begin{equation*}
 |f|_{\mathrm{sup}}=\sup_{x\in\mathrm{Spm}(\matheur{A})}|f(x)|.
\end{equation*}
It is known that for any representation 
$\alpha\colon K\{t_1,\ldots,t_n\}/I\stackrel{\sim}{\rightarrow}\matheur{A}$,
we have $|f|_{\mathrm{sup}}\leq |f|_{\alpha}$ for every $f\in\matheur{A}$
(cf.\ \cite[6.2.1]{na}).

\begin{appthm}[Maximal modulus principle]
For an affinoid algebra 
$\matheur{A}$ and an element $f\in\matheur{A}$, there exists a point
$x\in\mathrm{Spm}(\matheur{A})$ such that $|f(x)|=|f|_{\mathrm{sup}}$.
\end{appthm}

We refer \cite[6.2.1]{na} for the proof.

Given an affinoid algebra $\matheur{A}$ over $K$, we introduce the following
notation:
\begin{equation*}
 \begin{array}{lcl}
 \matheur{A}^{\circ}&=&\{f\in\matheur{A}\ |\ |f|_{\mathrm{sup}}\leq 1\},\\
 \matheur{A}^{\circ\circ}&=&\{f\in\matheur{A}\ |\ |f|_{\mathrm{sup}}<1\},\\
 \ovl{\matheur{A}}&=&\matheur{A}^{\circ}/\matheur{A}^{\circ\circ};
 \end{array}
\end{equation*}
$\matheur{A}^{\circ}$ is a $R$-subalgebra of $\matheur{A}$, and 
$\matheur{A}^{\circ\circ}$ is an ideal of it.
For instance, $K^{\circ}=R$, $K^{\circ\circ}=\mathfrak{m}$, and $\ovl{K}=k$.

\begin{appexa}\label{exa-red-tate}
We can immediately calculate:
$$
\begin{array}{lrl}
{\matheur{T}^n_K}^{\circ}&=&R\{t_1,\ldots,t_n\}\\
&\colon=&\left\{
\begin{array}{r|l}
\sum_{\nu_1,\ldots,\nu_n\geq 0}
a_{\nu_1,\ldots,\nu_n}t_1^{\nu_1}\cdots t_n^{\nu_n}&
|a_{\nu_1,\ldots,\nu_n}|\rightarrow 0\ \mathrm{for}\\
\in R[[t_1,\ldots,t_n]]
&\nu_1+\cdots+\nu_n\rightarrow\infty
\end{array}
\right\}\\
\ovl{\matheur{T}^n_K}&=&k[t_1,\ldots,t_n].
\end{array}
$$
Note that ${\matheur{T}^n_K}^{\circ}$ is the $(\pi)$-adic completion of
the polynomial ring $R[t_1,\ldots,t_n]$ over $R$.
It should be noticed that $\ovl{\matheur{T}^n_K}$ is a 
{\it polynomial ring} over $k$; indeed, each element in
$R\{t_1,\ldots,t_n\}$ lies, by the convergence condition, in 
${\matheur{T}^n_K}^{\circ\circ}$ modulo finitely many terms.
\end{appexa}

\begin{apppro}
Let $\matheur{A}$ be an affinoid algebra over $K$. Then,
\begin{enumerate}
% \renewcommand{\labelenumi}{\arabic{enumi}.}
 \item $\matheur{A}^{\circ}$ is a model of $\matheur{A}$$;$
       i.e.\ $\matheur{A}^{\circ}\otimes_RK\cong\matheur{A}$,
 \item $\matheur{A}^{\circ}$ is $(\pi)$-adically completed, 
 \item $\matheur{A}^{\circ}$ is topologically of finite
       type over $R$$;$ i.e.\
       $R\{t_1,\ldots,t_n\}/\mathfrak{a}\cong\matheur{A}$ 
       for some ideal $\mathfrak{a}$ in $R\{t_1,\ldots,t_n\}$,
 \item $\matheur{A}^{\circ}$ is flat over $R$,
 \item $\ovl{\matheur{A}}$ is a $k$-algebra of finite type.
\end{enumerate}
\end{apppro}

\begin{proof}
(i) and (ii) are easy to see.
The flatness of $R$ is now equivalent to the lack of $R$-torsion.
So (iv) can be seen immediately.

Before proving (iii) and (v), we need some terminology:
let $\|\cdot\|_{\alpha}$ be a Banach norm of $\matheur{A}$ ($\alpha\colon
\matheur{T}^n_K\rightarrow\matheur{A}$: surjective).
We say an element $f$ of $\matheur{A}$ is {\it power-bounded} if 
$\{\|f^n\|_{\alpha}\ |\ n\in\N\}$ is bounded.
This does not depend on the choice of the presentation $\alpha$ since
the equivalence class of the resulting norms do not change.
An element $f$ of $\matheur{A}$ is said to be {\it topologically nilpotent}
if $\lim f^n=0$; here the limit is taken with respect to the Banach
norm $\|\cdot\|_{\alpha}$, and is not dependent of the choice.

Let us prove (iii).
By \cite[6.2.3]{na}, $\matheur{A}^{\circ}$ is the set of all power-bounded
elements of $\matheur{A}$, and $\matheur{A}^{\circ\circ}$ is the set of all
topologically nilpotent elements.
Hence a choice of a presentation 
$\alpha\colon\matheur{T}^n_K\rightarrow\matheur{A}$ induces surjections
$$
\begin{array}{lcl}
{\matheur{T}^n_K}^{\circ}&\longrightarrow&\matheur{A}^{\circ}\\
{\matheur{T}^n_K}^{\circ\circ}&\longrightarrow&\matheur{A}^{\circ\circ}.
\end{array}
$$
(iii) is due to the surjectivity of the first arrow.
By the surjectivity of these arrows and \ref{exa-red-tate}, we have
(v).
\end{proof}

\begin{appdfn}\label{dfn-admissiblealg}
An $R$-algebra $A$ is said to be {\it admissible} if $A$ is 
$(\pi)$-adically completed, flat over $R$, and topologically 
of finite type over $R$.
\end{appdfn}

\begin{apprem}
By \cite[1.1\ (c)]{farg1rs}, any admissible $R$-algebra $A$ is 
topologically of finite presentation, i.e., there exists a finitely
generated ideal $\mathfrak{a}$ in $R\{t_1,\ldots,t_n\}$ such that
$R\{t_1,\ldots,t_n\}/\mathfrak{a}\cong\matheur{A}$.
\end{apprem}

\begin{appdfn}[Formal model and Analytic reduction]
Let $X=\mathrm{Spm}\matheur{A}$ be an affinoid over $K$.
Then the formal scheme $\mathscr{X}=\Spf\matheur{A}^{\circ}$ over $R$ is called
the {\it formal model} of $X$, and the algebraic scheme 
$\ovl{X}=\Spec\ovl{\matheur{A}}$ over $k$ is called the {\it analytic
reduction} of $X$.
\end{appdfn}

We define the so-called {\it reduction map}
$$
\mathrm{Red}_X\colon
X=\mathrm{Spm}\matheur{A}\longrightarrow\ovl{X}=\Spec\ovl{\matheur{A}}
$$
by $\mathfrak{a}\mapsto\mathfrak{a}\cap\matheur{A}^{\circ}/\mathfrak{a}\cap
\matheur{A}^{\circ\circ}$.

\begin{apppro}
The map $\mathrm{Red}_X$ is continuous with respect to the 
Zariski topology of $\ovl{X}$ and the G-topology of $X$; 
it maps $X$ surjectively onto the set of all
closed points of $\ovl{X}$.
\end{apppro}

\begin{proof}
Let us consider the affine open set 
$U_{\ovl{f}}=\Spec\ovl{\matheur{A}}_{\ovl{f}}$ of $\Spec\ovl{\matheur{A}}$
relevant to $\ovl{f}\in\ovl{\matheur{A}}$.
Take $f\in\matheur{A}^{\circ}$ from the residue class $\ovl{f}$.
The function $f$ takes values less than equal to $1$ on 
$\mathrm{Spm}\matheur{A}$.
So $U_f=\{x\in\mathrm{Spm}\matheur{A}\ |\ |f(x)|=1\}$ is a rational
subdomain whose associated affinoid is $\mathrm{Spm}A\{f,f^{-1}\}$.
Since the condition $|f(x)|=1$ is equivalent to 
$\ovl{f}(\mathrm{Red}_X(x))\neq 0$, we have 
$\mathrm{Red}_X^{-1}(U_{\ovl{f}})=U_f$.
Then the first assertion follows.
The last one is almost trivial.
\end{proof}

\begin{appexa}[Analytic reduction of the unit disk]
Let us consider the analytic reduction of the unit disk 
$\matheur{D}^1=\matheur{D}(0,1^{+})=\mathrm{Spm}K\{T\}$
(here we assume that $K$ is algebraically closed for simplicity).
The analytic reduction is given by the affine line 
$\mathbb{A}^1_k=\Spec k[T]$.
The reduction map $\mathrm{Red}_{\matheur{D}^1}$ maps all ``interior''
points, i.e.\ those points lying in $\matheur{D}(0,1^{-})$ to the origin
in $\mathbb{A}^1_k$, and the other points to the non-zero points.
If we identify $\matheur{D}^1$ and $R$, $\mathbb{A}^1_k$and $k$, then
this is nothing but the mapping obtained by the reduction modulo
$\mathfrak{m}$.
\end{appexa}

It is very important to see how the formal models behave under 
localization of affinoids.
Let $X=\mathrm{Spm}\matheur{A}$ be an affinoid over $K$ and $f_0,\ldots,f_n\in
\matheur{A}$ a sequence of elements which does not have common zeros
on $\mathrm{Spm}\matheur{A}$.
We consider the admissible covering $X=\bigcup_iU_i$ consisting of 
rational subdomains
$$
U_i=\{x\in X \mid |f_j(x)|\leq|f_i(x)|\ \mathrm{for}\ j\neq i\};
$$
by multiplying a suitable power of $\pi$, we may assume that each
$f_i$ belongs to $\matheur{A}^{\circ}$.
For each $i$, 
\begin{equation*}
 U_i=\mathrm{Spm}\matheur{A}\left\{\frac{f_0}{f_i},\ldots,
			     \frac{f_n}{f_i}\right\},
\end{equation*}
and then the corresponding formal model is given by 
\begin{equation*}
 \mathscr{U}_i=\Spf\matheur{A}^{\circ}
 \left\{\frac{f_0}{f_i},\ldots,\frac{f_n}{f_i}\right\}/
 \left((\pi)\mathrm{-torsions}\right).
\end{equation*}
These formal schemes glue and the resulting morphism 
\begin{equation*}
 \bigcup_{0\leq i\leq n}\mathscr{U}_i\longrightarrow\mathscr{X}
  =\Spf\matheur{A}^{\circ}
\end{equation*}
is the formal blow-up along the ideal 
$\mathfrak{a}=(f_0,\ldots,f_n)\subset\matheur{A}^{\circ}$.
We note here that the ideal $\mathfrak{a}$ contains a power of $\pi$, or
what amounts to the same, $\mathfrak{a}$ is an open ideal, for 
$\mathfrak{a}\matheur{A}=\matheur{A}$.

This simple observation indicates:

\vspace{1ex}\noindent
{\bf Slogan (vague)}: {\it Refinements of admissible coverings corresponds
to formal blow-up's along open coherent ideals.}

\vspace{1ex}

\begin{appexa}
Let us assume in this example that $K$ is algebraically closed.
We consider the admissible covering $\matheur{D}^1=U_1\cup U_2$ of the
unit disk $\matheur{D}^1=\mathrm{Spm}K\{T\}$ given by
\begin{equation*}
 \begin{array}{cclcl}
 U_1&=&\{z\in K\ |\ |z|\leq|\pi|\}&=&\mathrm{Spm}K\left\{\frac{T}{\pi}
 \right\},\\
 U_2&=&\{z\in K\ |\ |\pi|\leq |z|\leq 1\}&=&
 \mathrm{Spm}K\left\{T,\frac{\pi}{T}\right\}.
 \end{array}
\end{equation*}
The corresponding formal models are
\begin{equation*}
 \begin{array}{cclcl}
 \mathscr{U}_1&=&\Spf R\left\{\frac{T}{\pi}\right\}
 &=&\Spf R\left\{T,U\right\}/(\pi U-T),\\
 \mathscr{U}_2&=&\Spf R\left\{T,\frac{\pi}{T}\right\}
 &=&\Spf R\left\{T,V\right\}/(TV-\pi)
 \end{array}
\end{equation*}
respectively. We glue them as $U=V^{-1}$ and get 
$\mathscr{U}_1\cup\mathscr{U}_2\rightarrow\mathscr{D}^1=\Spf R\{T\}$,
which is the formal blow-up along the ideal $(\pi,T)$.
Meanwhile, the analytic reduction is obtained by gluing the following 
two affine sets:
\begin{equation*}
 \begin{array}{ccl}
 \ovl{U}_1&=&\Spec k\left[T\right],\\
 \ovl{U}_2&=&\Spec k\left[T,V\right]/(TV),
 \end{array}
\end{equation*}
this is the normal crossing of the affine line with coordinate $T$ and 
the projective line with the inhomogeneous coordinate $V$.
\end{appexa}

We are going to state the Raynaud's theorem, which provides a close
relationship between rigid geometry and formal geometry; we need some
terminology:

\begin{appdfn}[Admissible formal scheme]
An $R$-formal scheme\footnote{We follow 
\cite[\S 10]{ega1} for generalities of formal schemes.}
$\mathscr{X}$ is said to
be {\it admissible} if it is Zariski locally isomorphic to admissible
$\Spf A$ with $A$ admissible in the sense of Definition 
\ref{dfn-admissiblealg}.
\end{appdfn}

\begin{appdfn}[Admissible formal blow-up]
Let $\mathscr{X}$ be an admissible formal scheme.
An {\it admissible formal blow-up} of $\mathscr{X}$ is the formal blow-up
along a coherent open ideal $\mathscr{I}\subset\OO_{\mathscr{X}}$, i.e., 
\begin{equation*}
 \mathscr{X}'=\lim_{\stackrel{\longrightarrow}{\lambda}}\mathrm{Proj}
 \bigoplus^{\infty}_{n=0}\left(\mathscr{I}^n\otimes_{\mathscr{O}_{\mathscr{X}}}
 \OO_{\mathscr{X}}/\pi^{\lambda}\right)\longrightarrow\mathscr{X}.
\end{equation*}
\end{appdfn}

We are going to construct the functor from the category of admissible 
formal schemes to the category of rigid analytic spaces; let us start 
by observing the affine case.
From an affine admissible formal scheme $X=\Spf A$, it is easy to
get the corresponding rigid analytic space: we just set 
\begin{equation*}
 A_{\mathrm{rig}}\colon=A\otimes_RK
\end{equation*}
which is an affinoid algebra over $K$.
If $A$ is presented as $A=R\{t_1,\ldots,t_n\}/\mathfrak{a}$, then
$K\{t_1,\ldots,t_n\}/\mathfrak{a}K\{t_1,\ldots,t_n\}$ gives a presentation
of $A_{\mathrm{rig}}$.

To globalize, we need to see the compatibility with localization.
Let us consider the completed localization $A\{f^{-1}\}$ of $A$ with
respect to $f\in A$.
The corresponding affinoid algebra is
\begin{equation*}
 \begin{array}{ccccl}
 A\{f^{-1}\}\otimes_RK&=&A\{T\}/(1-Tf)\otimes_RK\\
 &=&A_{\mathrm{rig}}\{T\}/(1-Tf)&=&A_{\mathrm{rig}}\{f^{-1}\},
 \end{array}
\end{equation*}
and the latter is the affinoid algebra attached to the
rational subdomain
\begin{equation*}
 \{x\in\mathrm{Spm}A_{\mathrm{rig}}\ |\ |f|\geq 1\};
\end{equation*}
this means that the completed localization of formal schemes just 
corresponds to the localization of rigid spaces with respect to the
G-topology.
So we globalize the recipe as above to obtain the functor
\begin{equation*}
 \mathrm{Rig}\colon\mathscr{X}\mapsto\mathscr{X}_{\mathrm{rig}}
\end{equation*}
from admissible formal $R$-schemes to $K$-rigid analytic spaces.
The rigid analytic space $\mathscr{X}_{\mathrm{rig}}$ is called the
{\it Raynaud generic fiber} of $\mathscr{X}$.

\begin{apppro}\label{pro-blow-up}
The functor $\mathrm{Rig}$ maps admissible formal blow-up's of 
admissible formal schemes to isomorphisms of rigid analytic spaces.
\end{apppro}

\begin{proof}
This has already been essentially shown just before the Slogan.
Indeed, we may limit ourselves to the affine case, for the centers
of blow-up's are coherent.
We have seen in this case that the Raynaud generic fiber of 
admissible formal blow-up's is the decomposition into rational
subdomains, which does not change the analytic structure.
\end{proof}

\vspace{1ex}
Now we state the theorem of Raynaud:
\begin{appthm}[Raynaud\ 1972]
The functor gives the equivalence of the following categories,
\begin{enumerate}
\item the category of quasi-compact admissible formal $R$-schemes,
      localized by admissible formal blow-up's,
\item the category of quasi-compact and quasi-separated rigid analytic
      spaces over $K$.
\end{enumerate}
\end{appthm}

The reader finds a nice comprehensive account of the proof of this
theorem in \cite{farg1rs}.


\section{Topology of Rigid Analytic Space.}

In this section, we assume, for simplicity, that $K$ is algebraically
closed.

\begin{appdfn}
 \begin{enumerate}
  \renewcommand{\labelenumi}{(\arabic{enumi})}
  \item A rigid analytic space $X$ is said to be {\it quasi-compact} if
	every admissible covering of $X$ has a finite admissible 
	refinement.
  \item A rigid analytic space $X$ is said to be {\it connected} if there
	is no admissible covering $\{U_i\}_{i\in I}$ of $X$
	such that 
	\begin{equation*}
	 \bigcup_{i\in I_1}U_i\cap\bigcup_{i\in I_2}U_i=\emptyset\ 
	  \mathrm{and}\ 
	  \bigcup_{i\in I_1}U_i\neq\emptyset\neq\bigcup_{i\in I_2}U_i
	\end{equation*}	
	for some non-empty subsets $I_1,I_2\subseteq I$ with $I=I_1\sqcup I_2$.
 \end{enumerate}
\end{appdfn}

Since admissible coverings of affinoids are fixed as finite, 
every affinoid is quasi-compact.
So, we see that a rigid analytic space is quasi-compact if and only if it is
a finite union of affinoids.
It is easy to see that a rigid analytic space $X$ is connected if 
and only if the ring $\Gamma(X,\OO_X)$ has no other idempotent than 
$0$ and $1$.

\begin{appdfn}
 A morphism $\pi\colon Y\rightarrow X$ of rigid analytic spaces over $K$
 is said to be an {\it analytic covering} if there exists an admissible
 covering\footnote{or else a topological covering, referring to the
 Grothendieck topology; cf. \chapref{chap-ana-p-period},
 \subsecref{sub-top-cov-et-cov}} $\{X_i\}_{i\in I}$ of $X$ such that for
 each $i\in I$ $\pi^{-1}(X_i)$ is isomorphic to the disjoint union of
 copies of $X_i$.  A connect rigid analytic space $X$ is said to be {\it
 simply connected} if there is no other connected analytic covering over
 $X$ than the trivial one $\mathrm{id}\colon X\rightarrow X$ (cf.\
 \cite{anopu}).
\end{appdfn}

\begin{applem}\label{lem-simply-connected}
Let $X$ be a connect rigid analytic space $X$.

$\mathrm{(1)}$
The space $X$ is simply connected if and only if every locally
constant sheaf of sets, with respect to the Grothendieck topology, 
is constant.

$\mathrm{(2)}$
If there exists an admissible covering $\{X_i\}_{i\in\N}$ of $X$
such that $(a)$ every $X_i$ is simply connected, and $(b)$
$X_i\subseteq X_{i+1}$ for all $i$, then $X$ is simply connected.
\end{applem}

\begin{proof}
Both are straightforward.
\end{proof}

\begin{appexa}[Topology of the closed disk]\label{para-closed-disk}
Let $\matheur{D}=\matheur{D}(0,1^{+})$ be the $1$-dimensional closed disk of
radius $1$.
A rational subdomain in $\matheur{D}$ of form
\begin{equation*}
 \{x\in\matheur{D}\ |\ |x-a|\leq\rho\ \mathrm{and}\ |x-a_i|\geq\rho_i\ 
 \mathrm{for}\ i=1,\ldots,s\},
\end{equation*}
where $\rho,\rho_i\in |K^{\ast}|$ and $a,a_i\in\matheur{D}$, is called a
{\it standard domain}.
This is the complement of finite union of open disks.
The following properties are easy to verify:
\begin{itemize}
\item If $S_1$ and $S_2$ are standard domains such that $S_1\cap S_2
\neq\emptyset$, then so are $S_1\cup S_2$ and $S_1\cap S_2$.
\item Every finite union of standard domains is uniquely decomposed 
as a disjoint union of standard domains.
\end{itemize}

What is more interesting is the following proposition:
\end{appexa}

\begin{apppro}\label{pro-standart-domain}
Every rational subdomain in $\matheur{D}$ is a finite union of standard
domains.
$($Hence standard domains generate the Grothendieck topology of 
$\matheur{D}$.$)$
\end{apppro}

\noindent
\begin{proof}[Sketch of Proof] (cf.\ \cite[III.1.18]{sgamc}).\ 
Let $R$ be the rational subdomain given by $|f_i(x)|\leq |f_0(x)|$
for $i=1,\ldots,n$.
Deforming those functions slightly, if necessary, we may assume 
$f_i$ and $f_j$ have no common zero for any $i\neq j$, and thus,
we may concentrate to the single inequality $|f_1(x)|\leq |f_0(x)|$.
By Weierstrass Preparation Theorem we may assume $f_0$ and $f_1$ are
polynomials having all their roots in $\matheur{D}$.
Then the proposition follows from an easy calculation.
\end{proof}

\begin{appcor}\label{cor-projective-line}
The Grothendieck topology of the projective line $\mathbb{P}^{1,\mathrm{an}}_K$
is generated by standard domains.
\end{appcor}

\begin{apppro}\label{pro-standard}
Standard domains are connected.
Moreover, every connected rational subdomain in
$\mathbb{P}^{1,\mathrm{an}}_K$ is simply connected.
\end{apppro}

\begin{proof}
The first assertion is obvious since the corresponding affinoid algebra
\begin{equation*}
 K\{t,t_0,\ldots,t_s\}/((t-a)-t_0\pi,\pi_i-t_i(t-a_i)\ \mathrm{for}\ 
 i=1,\ldots,s),
\end{equation*}
where $a,a_i\in K$ and $\pi,\pi_i\in K$, is integral.
(Note that $K\{t,t_0,\ldots,t_s\}$ is Noetherian and factorial.)
The second one is proved by using \ref{lem-simply-connected} (1).
Let $S$ be a standard domain and $\mathscr{F}$ a constant sheaf of
sets on $S$.
We can take an admissible covering $\{S_i\}$ of $S$ which trivializes 
$\mathscr{F}$.
The point is that any non-disjoint union of standard domains is 
standard, and hence, connected.
So the restriction maps $\mathscr{F}(S_i)\rightarrow
\mathscr{F}(S_i\cap S_j)$ must be bijective as far as $S_i\cap S_j
\neq\emptyset$.
It follows that $\mathscr{F}(S_i\cup S_j)
\rightarrow\mathscr{F}(S_i)$ are bijective, and repeating this argument,
we concludes $\mathscr{F}$ is constant.
\end{proof}

\vspace{1ex}
The following corollary may seen at first:

\begin{appcor}\label{cor-simply-connected}
Every connected open subanalytic space\footnote{We refer \cite[9.3.1]{na} 
for the definition of open subanalytic spaces.} in 
$\mathbb{P}^{1,\mathrm{an}}_K$ is simply connected.
\end{appcor}

\begin{proof}
A connected open subanalytic space $U$ has an admissible covering 
consisting of standard domains.
By \cite[III,2.6]{sgamc} and the easy fact that a non-disjoint union of 
connected sets is again connected, we see that a finite non-disjoint
union of standard domain is a connected rational subdomain.
Hence the corollary follows from Lemma \ref{lem-simply-connected} (2)
and Proposition \ref{pro-standard}.
\end{proof}

\vspace{1ex}
Finally we quote, without proof, a theorem by van der Put 
which indicates a close relation between the topology 
of a rigid analytic space and that of its analytic reduction:

\begin{appthm}\label{thm-vanderput}
Let $X$ be a quasi-compact rigid analytic space over $K$ which has 
the irreducible and smooth analytic reduction. Then $X$ is simply
connected.
\end{appthm}

For the proof, see \cite{anopu}.
\begin{apprem}
By the corollary, we know, for example, that the space
$\mathbb{G}^{\mathrm{an}}_K=K^{\times}$ is simply connected.  But, on
the other hand, it is easy to see that the map $K^{\times} \rightarrow
K^{\times}$ by $x\mapsto x^n$ is an analytic morphism.  As far as the
characteristic of $K$ does not divide $n$, we are tempted to think of
this morphisms also as an analytic covering.  But it is not.  Certainly,
this morphism induces isomorphisms between stalks of the structure
sheaf, but there is no {\it admissible} covering which trivializes this
morphism! 
This kind of morphisms is said to be {\it \'etale}; analytic coverings
are \'etale, but not vice versa.  Simply connected analytic spaces
may have many \'etale coverings.
\end{apprem}

\section{Berkovich' approach to non-archimedean analysis}

Berkovich' viewpoint provides an innovative approach to topological
problems of $p$-adic geometry and analysis, which is much closer to the
familiar intuition derived from the complex context
\cite{staagonf}\footnote{Already in \cite{coas} the {\it esquisse} of
Berkovich' idea can be seen.}.  His analysis, in fact, generalizes, or
properly speaking, ``completes'' the rigid analysis.  This situation is,
according to what he says in the introduction of his book, somehow
analogous to that of $\R$ completing $\Q$.  The standard archimedean
metric space $\Q$ is totally disconnected.  The rigid analytic viewpoint
corresponds, to some extent, to regarding every rational interval
$\{r\in\Q\ |\ a\leq r\leq b\}$ as ``connected'', and considering only those
functions which come from analytic functions on the corresponding real
interval.  Berkovich' non-archimedean analysis can be compared, in this
context, to analysis on $\R$ itself. Namely, his presentation of
analytic spaces complements those of rigid analysis by putting some
``generic points'', and consequently, simplifies their topology.

Let us be a little more precise about his theory.
Let $\matheur{A}$ be a Banach ring, i.e., a ring together with a norm 
$\|\cdot\|$ with respect to which $\matheur{A}$ is complete.

\begin{appdfn}\label{dfn-seminorm}

(1) A semi-norm $|\cdot |$ on $\matheur{A}$ is said to be {\it bounded} 
if there exists $C>0$ such that $|f|\leq C\|f\|$ for any $f\in\matheur{A}$.

(2) A semi-norm $|\cdot |$ on $\matheur{A}$ is said to be 
{\it multiplicative} if $|fg|=|f||g|$ for all $f,g\in\matheur{A}$.
\end{appdfn}

Replacing the maximal spectra $\Spm(\matheur{A})$ of rigid geometry, we have:
\begin{appdfn}\label{dfn-Berkovich-spectrum}
The Berkovich' {\it spectrum} $\matheur{M}(\matheur{A})$ for a Banach ring
$\matheur{A}$ is the set of all bounded multiplicative semi-norms on
$\matheur{A}$, endowed with the weakest topology so that the real
valued functions of form $|\cdot |\mapsto |f|$ for $f\in\matheur{A}$
are continuous.
\end{appdfn}

In case $\matheur{A}$ is an affinoid algebra (a {\it strict} affinoid
algebra, in Berkovich' term) over $K$, the affinoid space 
$\mathrm{Spm}(\matheur{A})$ is naturally viewed as a subspace, with 
the relative topology, of the Berkovich' spectrum $\matheur{M}(\matheur{A})$.
This is done by identifying an element $\mathfrak{m}$ in
$\mathrm{Spm}(\matheur{A})$ with the unique semi-norm $|\cdot |$ on 
$\matheur{A}$ such that $\mathfrak{m}=\Ker |\cdot |$.
The space $\mathrm{Spm}(\matheur{A})$ is identified, in terms of this 
correspondence, with the subspace of $\matheur{M}(\matheur{A})$ consisting of
semi-norms $|\cdot |$ such that $\dim_K\matheur{A}/\Ker |\cdot |$ is 
finite (the ``classical points'').

\begin{appexa}
Let us consider the usual absolute-value norm $|\cdot|_{\infty}$
on $\Z$. The spectrum
$\matheur{M}(\Z)$ consists of the following points:
\begin{enumerate}
\item $|\cdot|_{\infty,\epsilon}\colon=|\cdot|_{\infty}^{\epsilon}$
($0<\epsilon\leq 1$).
\item The $p$-adic norm $|\cdot|_{p,\epsilon}\colon$ with 
$|p|_{p,\epsilon}=\epsilon$ ($0<\epsilon\leq 1$).
\item The semi-norm $|\cdot|_p$ induced from the trivial norm on $\Z/p\Z$.
\item The trivial norm $|\cdot|_0$.
\end{enumerate}
As a topological space, $\matheur{M}(\Z)$ is a tree with end points
$|\cdot|_p$ and $|\cdot|_{\infty}$. Each of these points is connected
by a single edge with $|\cdot|_0$.
\end{appexa}

On the Berkovich' spectrum, the {\it affinoid subdomains} (incl.\ {\it 
rational subdomains}) are defined in the similar way to rigid 
analysis. 
Thus one can define the structure sheaves on spectra, which 
one calls affinoid spaces.
The local data consisting of Berkovich' affinoids and the structure
sheaves can glue to locally ringed spaces; thus we obtain {\it analytic
spaces} in Berkovich' sense.

\begin{appexa}
(Cf.\ \cite[1.5]{staagonf}.)
Let $K$ be a non archimedean complete field. An affine space over $K$ is
defined by
\begin{equation*}
 \mathbb{A}^n_K=\{\textrm{multiplicative seminorm}\ |\cdot|\ \textrm{on}\ 
 K[t_1,\ldots,t_n]\ |\ |\cdot|_{|_K}\ \textrm{is bounded}\}
\end{equation*}
together with the weakest topology so that $|\cdot|\mapsto |f|$
($f\in K[t_1,\ldots,t_n]$) is continuous.
This has a structure of analytic space so that the analytic functions
over it is characterized by limits of rational functions.
When $K=\C$ the space $\mathbb{A}^n_K$ is nothing but the usual affine
space in the complex analytic sense.
If $K$ is a non-archimedean field, then $\mathbb{A}^n_K$ is a 
union of balls; i.e., the spectrum of Tate algebra.
\end{appexa}

Here are some general properties of Berkovich' analytic spaces:
\begin{enumerate}
\item Every connected $K$-analytic space is arcwise connected.
\item If $X=\matheur{M}(\matheur{A})$ with $\matheur{A}$ an affinoid algebra over
$K$, then $\textrm{Krull-dim}\matheur{A}=\mathrm{dim}|X|$, where $|X|$
is the underlying topological space of $X$.
\item If a $K$-analytic space $X$ is smooth, then it is locally
contractible.
\end{enumerate}

Relation between Berkovich' analytic spaces and rigid analytic spaces is
as follows.
If a Berkovich' $K$-analytic space $X$ is formed only by spectra of 
(strict) affinoid algebras, we say $X$ is strict.
There exists a functor from the category of separated strict 
$K$-analytic spaces to that of separated rigid analytic spaces over $K$
by 
\begin{equation*}
 X\mapsto X_0=\{x\in X\ |\ [K(x)\colon K]<\infty\},
\end{equation*}
where $K(x)$ is the residue field at $x$.

\begin{appthm}
This functor is fully-faithful, and preserves fiber products.
\end{appthm}
